L'esperienza descritta ha come scopo stabilire il valore della velocità della luce, grandezza fondamentale per lo studio della fisica sulla cui misura sono state fatte nei secoli numerose analisi. La genialità del metodo proposto da Foucault, e che è stato ricreato abbastanza fedelmente in laboratorio, sta nel fatto che non si cerca un valore della velocità in maniera classica (ovvero come rapporto tra una distanza percorsa e il tempo necessario per percorrerla, che richiederebbe o spazi enormi o capacità di percepire tempi molto piccoli), ma, sfruttando fenomeni ottici, si riesce a lavorare con grandezze più facilmente misurabili senza strumenti estremamente sofisticati.
L'esperienza si basa su un raggio laser che, colpendo uno specchio piano ruotante , viene indirizzato verso uno specchio concavo, che lo costringe a ritornare indietro seguendo il medesimo percorso d'andata. Nel frattempo, lo specchio rotante cambia posizione: il laser di ritorno, colpendo lo specchio rotante, non verrà mandato verso la sorgente, ma verrà leggermente deviato. Conoscendo la dimensione di questa deviazione e alcune costanti proprie dell'apparato sperimentale (come la velocità angolare di rotazione dello specchio), risulta possibile calcolare la velocità della luce.

Purtroppo, non è stato possibile eseguire l'esperienza nel caso di rotazione in senso orario dello specchio a causa di un malfunzionamento di uno strumento fondamentale: il sensore in grado di misurare e comunicare la frequenza di rotazione dello specchio. \'E risultato infatti che tale strumento non era in grado di funzionare in maniera esatta quando si superava una frequenza di 400 Hz, rendendo impossibile ottenere misure valide della velocità della luce.

Questo problema non si è manifestato invece per le rotazioni in senso antiorario, dove si sono registrati numerosi dati che hanno permesso di calcolare la velocità della luce $c$ come:
\begin{equation}
    c_{CCW}=(2.997\pm0.078)\times 10^8\frac{m}{s}.
\end{equation}
Questo valore è stato confrontato con la velocità reale $c=(2.9979\times 10^8)\frac{m}{s}$ mediante il confronto tramite t di Student, e si è registrato 98\% di probabilità che la differenza abbia avuto origine statistica.

\'E stato poi possibile sfruttare una funzionalità del motorino collegato allo specchio rotante che permetteva di raggiungere la massima frequenza in entrambe le direzioni di rotazione e confrontare la posizione del raggio laser di ritorno nei due casi. Grazie a tale set di dati si è trovato un secondo valore della velocità della luce, pari a 
\begin{equation}
    c_{turbo}=(3.007\pm0.043)\times 10^8 \frac{m}{s}.
\end{equation}
Confrontando nuovamente il valore ottenuto con il valore atteso mediante il test del t di Student, risulta che esiste l'84\% si probabilità che la differenza tra i due abbia origini statistiche.

In entrambi i casi, sia nella rotazione antioraria che nell'utilizzo del "turbo", si sono ottenuti valori statisticamente coerenti al valore atteso.


